% Replace the response after the existing AE comment; do not duplicate its label.
% Requires xcolor, natbib, and entries from AE1_references.bib.
% The closing paragraph refers only to the two additions the author confirmed.
{\color{blue}
\noindent\textbf{Response:}
We thank the Associate Editor for suggesting this perspective. It helps clarify how LEAN-LLM-OPT combines reference modeling knowledge with reusable procedural guidance, and how this design relates to agent skills.

Prior work distinguishes declarative knowledge, concerning facts and concepts, from procedural knowledge, concerning how to perform a task \citep{ae1_li2024metacognitive}. Retrieval-augmented generation (RAG), in turn, is a mechanism for conditioning generation on retrieved external material \citep{ae1_lewis2020rag}. Retrieval describes how information is accessed, rather than whether its content is declarative or procedural. For example, Buffer of Thoughts retrieves reusable reasoning templates and instantiates them for new problems \citep{ae1_yang2024buffer}. Thus, retrieved material can supply modeling knowledge or procedural guidance; retrieval alone does not make that material a skill.

Skills organize procedural knowledge for reuse and can incorporate reference knowledge. \citet{ae1_zhou2026skillsurvey} describe skills in terms of instructions, supporting resources, and applicability conditions. OpenAI's official documentation similarly describes a skill as a package of instructions and supporting files, including reference material \citep{ae1_openai2026skills}. It also describes skills that coordinate search and fetch tools and apply reusable procedures to the information returned \citep{ae1_openai2026skilltools}. These descriptions support combining reference material and procedural instructions within a skill, as well as retrieving relevant material during its execution. They do not equate retrieved content with a skill, and a skill may obtain additional task-relevant information at runtime.

In LEAN-LLM-OPT, Ref-Data primarily supplies declarative modeling knowledge through problem descriptions, associated data, expert-provided formulations, and annotations identifying relevant data. The workflow generation agent integrates selected reference content with instructions for data retrieval and model construction. In the type-tailored workflow, the reference query, relevant-data category, data details, and formulation populate the Question, Action Input, Observation, and Final Answer components. Type-aware formatting and a type-specific output structure show how data enter the model. The resulting demonstration therefore combines modeling content with a procedure for using it. The model generation agent follows and adapts this guidance to the target query and datasets. For Others and Mixture, the type-agnostic workflow instead guides modeling through an abstract model plan based on the query and dataset schema.

Viewed functionally, these modeling procedures are \emph{skill-like}: classification identifies the applicable route, the instructions guide its execution, and the reference cases provide supporting knowledge. The connection to skills therefore lies in the organized use of reference knowledge, not in reference retrieval alone. Knowledge and procedural guidance are already combined within the constructed demonstration; their roles remain distinguishable within that representation.

We distinguish this connection from particular packaging and learning mechanisms. LEAN-LLM-OPT implements its guidance through prompts and workflow construction, rather than the standalone skill packages evaluated in SkillsBench \citep{ae1_li2026skillsbench}. Its curated resources and reusable instructions remain fixed during evaluation, and completed trajectories are not used to acquire or revise a persistent skill library. Such learning is studied in Agent Workflow Memory \citep{ae1_wang2025awm} and, for optimization modeling, OptSkills \citep{ae1_yang2026optskills}. The absence of autonomous skill learning does not preclude skill-like procedural reuse; it specifies the scope of our implementation. Accordingly, LEAN-LLM-OPT remains an \emph{agentic workflow construction framework} that integrates reference modeling knowledge with reusable procedures for application to new instances. This inference-time guidance can also complement a fine-tuned model.

We have added two discussions to the manuscript. \emph{Knowledge and Skills in LLM Agents} in Section~1.2 (\emph{Related Works}) introduces the relevant terminology and literature. \emph{Connecting Reference Knowledge and Modeling Skills} in Section~3.3 (\emph{Novelties in Our Design and Key Messages}) explains how reference cases and procedural guidance work together, clarifies the roles of the workflow generation and model generation agents, and delineates the framework's scope relative to skill packaging and learning.
}
